\documentclass[10pt, aspectratio=169]{beamer}
% Preámbulo modular para presentaciones Beamer (Modelación Estadística)
% Diseño y paleta institucional del Tecnológico de Monterrey

\documentclass[aspectratio=169,xcolor={dvipsnames,table}]{beamer}

\usepackage[utf8]{inputenc}
\usepackage[spanish,mexico]{babel}
\usepackage{amsmath,amssymb,amsfonts,mathtools}
\usepackage{graphicx}
\usepackage{listings}
\usepackage{booktabs}
\usepackage{tikz}
\usetikzlibrary{matrix,arrows,decorations.pathmorphing,shapes.geometric,positioning}

% Paleta Institucional Tecnológico de Monterrey
\definecolor{TecAzul}{HTML}{0039A6}       % Azul TEC: color primario institucional
\definecolor{TecAzulOscuro}{HTML}{1A2E51} % Azul marino oscuro
\definecolor{TecRojo}{HTML}{EC2661}       % Rojo de acento (paleta miTec)
\definecolor{TecGris}{HTML}{646464}       % Gris neutro para comentarios y subtítulos
\definecolor{TecFondoBloque}{HTML}{F4F6F9}% Fondo suave para bloques

% Configuración de Tema Beamer
\usetheme{Boadilla}
\usecolortheme[named=TecAzulOscuro]{structure}
\setbeamercolor{alerted text}{fg=TecRojo}
\setbeamercolor{palette primary}{bg=TecAzulOscuro,fg=white}
\setbeamercolor{palette secondary}{bg=TecAzul,fg=white}
\setbeamercolor{palette tertiary}{bg=TecAzulOscuro!80!black,fg=white}
\setbeamercolor{title}{fg=TecAzulOscuro,bg=TecFondoBloque}
\setbeamercolor{frametitle}{fg=TecAzulOscuro,bg=TecFondoBloque}

% Bloques personalizados elegantes
\setbeamercolor{block title}{bg=TecAzulOscuro,fg=white}
\setbeamercolor{block body}{bg=TecFondoBloque,fg=black}
\setbeamercolor{block title alerted}{bg=TecRojo,fg=white}
\setbeamercolor{block body alerted}{bg=TecRojo!8,fg=black}
\setbeamercolor{block title example}{bg=TecAzul,fg=white}
\setbeamercolor{block body example}{bg=TecAzul!8,fg=black}

% Redondear bordes y añadir sombra sutil
\setbeamertemplate{blocks}[rounded][shadow=true]
\setbeamertemplate{navigation symbols}{}

% Pie de página limpio con número de diapositiva
\setbeamertemplate{footline}{
  \leavevmode%
  \hbox{%
  \begin{beamercolorbox}[wd=.7\paperwidth,ht=2.25ex,dp=1ex,left]{author in head/foot}%
    \hspace*{2ex}\insertshortauthor~~(\insertshortinstitute) \hfill \insertshorttitle
  \end{beamercolorbox}%
  \begin{beamercolorbox}[wd=.3\paperwidth,ht=2.25ex,dp=1ex,right]{date in head/foot}%
    \insertshortdate{}\hspace*{2em}
    \insertframenumber{} / \inserttotalframenumber\hspace*{2ex}
  \end{beamercolorbox}}%
  \vskip0pt%
}

% Configuración de Listings de Python para visualización pedagógica de código
\lstset{
    language=Python,
    basicstyle=\ttfamily\footnotesize,
    keywordstyle=\color{TecAzulOscuro}\bfseries,
    stringstyle=\color{TecRojo},
    commentstyle=\color{TecGris}\itshape,
    showstringspaces=false,
    breaklines=true,
    frame=lines,
    rulecolor=\color{TecAzulOscuro!30},
    backgroundcolor=\color{TecFondoBloque},
    aboveskip=0.5em,
    belowskip=0.5em,
    literate={á}{{\'a}}1 {é}{{\'e}}1 {í}{{\'i}}1 {ó}{{\'o}}1 {ú}{{\'u}}1
             {Á}{{\'A}}1 {É}{{\'E}}1 {Í}{{\'I}}1 {Ó}{{\'O}}1 {Ú}{{\'U}}1
             {ñ}{{\~n}}1 {Ñ}{{\~N}}1 {¿}{{?`}}1 {¡}{{!`}}1
}

% Metadatos institucionales por defecto
\title[Modelación Estadística]{Modelación Estadística para Ciencia de Datos e Ingeniería}
\author[J. Castillo Colmenares]{Juliho David Castillo Colmenares}
\institute[Tec de Monterrey]{Tecnológico de Monterrey \\ Escuela de Ingeniería y Ciencias}
\date{\today}

% Comandos y macros estadísticas compartidas para las presentaciones Beamer (Metropolis)

% Conjuntos numéricos básicos
\newcommand{\R}{\mathbb{R}}
\newcommand{\N}{\mathbb{N}}
\newcommand{\Q}{\mathbb{Q}}
\newcommand{\Z}{\mathbb{Z}}
\newcommand{\set}[1]{\left\{ #1 \right\}}
\newcommand{\abs}[1]{\left|#1\right|}
\newcommand{\norm}[1]{\left\| #1 \right\|}

% Comandos estadísticos y probabilísticos del curso
\newcommand{\Var}{\operatorname{Var}}
\newcommand{\cov}{\operatorname{Cov}}
\newcommand{\corr}{\rho}
\newcommand{\comb}[2]{\begin{pmatrix} #1 \\ #2 \end{pmatrix}}
\newcommand{\s}{\sigma}
\newcommand{\particion}{\mathfrak{P}}
\newcommand{\card}[1]{\#\left( #1 \right)}
\newcommand{\seq}[3]{\set{#1}_{#2}^{#3}}
\newcommand{\E}{\mathbb{E}}
\newcommand{\Prob}{\mathbb{P}}

% Entornos pedagógicos y cajas en español académico natural
\newenvironment{discusion}{%
  \begin{alertblock}{Pregunta de Discusión}%
}{%
  \end{alertblock}%
}

\newenvironment{reflexion}{%
  \begin{alertblock}{Reflexión para el Aula}%
}{%
  \end{alertblock}%
}

\newenvironment{paradoja}{%
  \begin{alertblock}{Paradoja de Exploración}%
}{%
  \end{alertblock}%
}

\newenvironment{motivacion}{%
  \begin{exampleblock}{Motivación: Ciencia de Datos en la Práctica}%
}{%
  \end{exampleblock}%
}

\newenvironment{retoclase}{%
  \begin{block}{Reto Rápido en Equipo}%
}{%
  \end{block}%
}


\title[Distribución Uniforme]{Sección 04.04: Distribución Uniforme Continua}
\subtitle{Fundamentos de $U(a, b)$, Cuantiles, Máxima Entropía y Orden Estadísticas}
\author[J. Castillo Colmenares]{Juliho Castillo Colmenares}
\institute{Tecnológico de Monterrey}
\date{\vspace{-1.2cm}}

\begin{document}

% Slide 1: Portada
\begin{frame}[plain]
  \titlepage
\end{frame}

% Slide 2: Hoja de Ruta
\begin{frame}{Hoja de Ruta --- Capítulo 04: Variables Aleatorias Continuas}
  \footnotesize
  ¿Dónde nos encontramos en el desarrollo modular del capítulo? \pause
  \begin{itemize}\itemsep=0.08em
    \item \textbf{04.01} Función de Densidad (PDF) y Soporte Continuo \pause
    \item \textbf{04.02} Función de Distribución Acumulada (CDF) Continua y Cuantiles \pause
    \item \textbf{04.03} Esperanza, Varianza y LOTUS Continuo \pause
    \item \textbf{04.04} Distribución Uniforme Continua $U(a, b)$ \emph{(Hoy)} \pause
    \item \textbf{04.05} Distribución Normal y Puntaje $Z$ \pause
    \item \textbf{04.06} Distribución Exponencial y Procesos sin Memoria \pause
    \item \textbf{04.07} Distribuciones Gamma, Beta y Weibull
  \end{itemize}
  \vspace{-0.05cm}
  \begin{block}{Objetivo de la Sesión}
    Introducir la distribución Uniforme Continua $U(a, b)$ como la más simple de las distribuciones continuas con soporte acotado, calcular sus momentos, demostrar que maximiza la entropía entre distribuciones con soporte dado, y aplicarla al método de inversión y a la prueba de Kolmogorov-Smirnov.
  \end{block}
\end{frame}

% Slide 3: Motivación
\begin{frame}{La Distribución Más Simple: Equiprobabilidad Continua}
  \footnotesize
  \begin{motivacion}[¿Por qué estudiar la uniforme?]
    La distribución Uniforme Continua $U(a, b)$ modela el caso más fundamental: una variable que toma cualquier valor en $[a, b]$ con la misma \emph{densidad}. Es la base de la simulación Monte Carlo (todas las muestras i.i.d. de cualquier distribución se generan a partir de uniformes), del análisis de inspección aleatoria, y de la prueba no paramétrica de Kolmogorov-Smirnov.
  \end{motivacion}

  \vspace{0.15cm}
  \pause
  \begin{itemize}\itemsep=0.1cm
    \item \textbf{Equiprobabilidad:} $f(x) = 1/(b-a)$ es constante en $[a, b]$. \pause
    \item \textbf{Máxima entropía:} $U(a, b)$ maximiza la incertidumbre entre distribuciones con soporte en $[a, b]$. \pause
    \item \textbf{Generador universal:} El método de inversión $X = F^{-1}(U)$ con $U \sim U(0, 1)$ produce muestras de cualquier distribución continua.
  \end{itemize}
\end{frame}

% Slide 4: Definición Formal
\begin{frame}{Definición Formal de la Distribución Uniforme $U(a, b)$}
  \footnotesize
  Una variable aleatoria continua $X$ tiene distribución \emph{uniforme continua} en $[a, b]$ con $a < b$ si su PDF es: \pause

  \vspace{0.1cm}
  \begin{block}{Función de Densidad de Probabilidad}
    $$f(x) = \begin{cases} \dfrac{1}{b-a} & a \le x \le b \\ 0 & \text{en otro caso} \end{cases}$$
    Por normalización: $\int_{a}^{b} \frac{1}{b-a}\,dx = 1$. En este caso escribimos $X \sim U(a, b)$.
  \end{block}

  \vspace{0.1cm}
  \pause
  \begin{alertblock}{CDF Lineal}
    \begin{align*}
      F(x) = P(X \le x) = \begin{cases} 0 & x < a \\ \dfrac{x-a}{b-a} & a \le x \le b \\ 1 & x > b \end{cases}
    \end{align*}
    La CDF es lineal con pendiente $1/(b-a)$ en el soporte. La inversa es trivial: $F^{-1}(p) = a + p(b-a)$.
  \end{alertblock}
\end{frame}

% Slide 5: Momentos
\begin{frame}{Esperanza y Varianza: $\E[X] = (a+b)/2$, $\Var(X) = (b-a)^{2}/12$}
  \footnotesize
  Los momentos de la distribución uniforme tienen formas cerradas simples: \pause

  \vspace{0.1cm}
  \begin{block}{Esperanza y Varianza}
    \begin{align*}
      \E[X] &= \int_{a}^{b} \frac{x}{b-a}\,dx = \frac{a + b}{2}, \\
      \Var(X) &= \int_{a}^{b} \frac{(x - (a+b)/2)^{2}}{b-a}\,dx = \frac{(b-a)^{2}}{12}, \\
      \sigma &= \frac{b-a}{\sqrt{12}} = \frac{b-a}{2\sqrt{3}}.
    \end{align*}
  \end{block}

  \vspace{0.1cm}
  \pause
  \begin{alertblock}{Propiedades Fundamentales}
    \begin{itemize}\itemsep=0.05cm
      \item \textbf{Simetría:} Si $X \sim U(a, b)$, entonces $Y = a + b - X \sim U(a, b)$ (simetría alrededor de la mediana).
      \item \textbf{Linealidad:} $Y = cX + d \sim U(ca+d, cb+d)$ para $c > 0$.
      \item \textbf{Cuantiles:} $q_{p} = a + p(b-a)$ son linealmente interpolados entre $a$ y $b$.
    \end{itemize}
  \end{alertblock}
\end{frame}

% Slide 6: Propiedad de Falta de Memoria
\begin{frame}{La Uniforme NO Tiene Propiedad de Falta de Memoria}
  \footnotesize
  A diferencia de la Exponential, la Uniforme \emph{no} tiene la propiedad de falta de memoria: \pause

  \vspace{0.1cm}
  \begin{block}{Diferencia con la Exponencial}
    Para $X \sim U(0, 15)$ y $a = 10$:
    \begin{align*}
      P(X \ge a + b \mid X \ge a) = \dfrac{P(X \ge a + b)}{P(X \ge a)} = \dfrac{(15 - a - b)^{+}}{(15 - a)} \ne e^{-b/15} = P(X \ge b).
    \end{align*}
    En general, la distribución Uniforme \emph{decrece} en probabilidad condicional a medida que $X$ se acerca al borde superior $b$.
  \end{block}

  \vspace{0.1cm}
  \pause
  \begin{alertblock}{Implicaciones Prácticas}
    \begin{itemize}\itemsep=0.05cm
      \item \textbf{Inspección aleatoria:} Si se sabe que $X \ge 10$ en $U(0, 15)$, la distribución condicional es $U(10, 15)$.
      \item \textbf{Aplicación en simulación:} Generar $U \sim U(0, 1)$ y luego $X = a + (b-a)U$ produce $X \sim U(a, b)$.
      \item \textbf{Pruebas de uniformidad:} La no-falta-de-memoria es un test diagnóstico: si los datos la exhiben, no son uniformes.
    \end{itemize}
  \end{alertblock}
\end{frame}

% Slide 7: Máxima Entropía
\begin{frame}{La Uniforme Maximiza la Entropía: $h(X) = \log(b-a)$}
  \footnotesize
  La distribución uniforme tiene la \emph{máxima entropía diferencial} entre todas las distribuciones con soporte en $[a, b]$: \pause

  \vspace{0.1cm}
  \begin{block}{Demostración por Multiplicadores de Lagrange}
    Maximizar $h(X) = -\int_{a}^{b} f(x) \log f(x)\,dx$ sujeto a $\int f = 1$. El Lagrangian es:
    \begin{align*}
      \mathcal{L} = -\int_{a}^{b} f \log f\,dx - \lambda \left( \int_{a}^{b} f\,dx - 1 \right).
    \end{align*}
    Derivando funcionalmente: $-\log f - 1 - \lambda = 0$, así $f(x) = e^{-(1+\lambda)}$ es constante. La normalización fuerza $f(x) = 1/(b-a)$.
  \end{block}

  \vspace{0.1cm}
  \pause
  \begin{alertblock}{Interpretación y Aplicaciones}
    \begin{itemize}\itemsep=0.05cm
      \item \textbf{Entropía uniforme:} $h(U(a,b)) = \log(b-a)$ nats. Uniform(0, 1) tiene entropía 0.
      \item \textbf{Inferencia bayesiana:} Como prior no informativo, la uniforme expresa "ignorancia total" sobre el parámetro.
      \item \textbf{Comparación:} Beta(2, 2) tiene entropía $\approx -0.125 < 0 = h(U(0,1))$, confirmando la maximalidad.
    \end{itemize}
  \end{alertblock}
\end{frame}

% Slide 8: Método de Inversión
\begin{frame}{Método de Inversión: $X = F^{-1}(U)$ con $U \sim U(0, 1)$}
  \footnotesize
  La uniforme es la base del \emph{método de inversión} para generar muestras de cualquier distribución continua: \pause

  \vspace{0.1cm}
  \begin{block}{Teorema de Inversión}
    Si $U \sim U(0, 1)$ y $F$ es una CDF continua estrictamente creciente, entonces $X = F^{-1}(U)$ tiene CDF $F$.
    \begin{proof}
      $P(X \le x) = P(F^{-1}(U) \le x) = P(U \le F(x)) = F(x)$ por monotonía de $F^{-1}$ y uniformidad de $U$.
    \end{proof}
  \end{block}

  \vspace{0.1cm}
  \pause
  \begin{alertblock}{Ejemplos de Generación}
    \begin{itemize}\itemsep=0.05cm
      \item \textbf{Exponential($\lambda$):} $X = -\frac{1}{\lambda}\ln(1-U) \approx -\frac{1}{\lambda}\ln U$.
      \item \textbf{Cauchy:} $X = \tan(\pi(U - 1/2))$.
      \item \textbf{Log-normal:} $X = \exp(\mu + \sigma Z)$ con $Z = \Phi^{-1}(U)$.
    \end{itemize}
  \end{alertblock}
\end{frame}

% Slide 9: Aplicaciones
\begin{frame}{Aplicaciones en Simulación, Inferencia y Ciencia de Datos}
  \footnotesize
  La Uniforme es la distribución más usada en simulación e inferencia: \pause

  \vspace{0.15cm}
  \begin{itemize}\itemsep=0.12cm
    \item \textbf{Simulación Monte Carlo:} Generación de muestras i.i.d. de cualquier distribución vía el método de inversión. \pause
    \item \textbf{Muestreo Aleatorio Simple:} La uniforme es la base del muestreo probabilístico en encuestas y experimentos. \pause
    \item \textbf{Prueba de Kolmogorov-Smirnov:} Compara la CDF empírica con una CDF teórica (uniforme o no). \pause
    \item \textbf{Generación de Números Pseudoaleatorios:} Algoritmos congruenciales lineales producen aproximaciones de $U(0, 1)$.
  \end{itemize}
\end{frame}

% Slide 12A-1: Lab Python (Bloque 1)
\begin{frame}[fragile]{Laboratorio en Python: PDF y CDF Uniforme}
  \scriptsize
  Validación de $\int f = 1$, comparación de CDF manual vs SciPy, y cuantiles (\texttt{04.04\_uniform\_distribution.py}):
  \vspace{0.02cm}
  {\lstset{basicstyle=\fontsize{4.5pt}{5.5pt}\selectfont\ttfamily}
  \lstinputlisting[firstline=15, lastline=40]{../../code/04_variables_aleatorias_continuas/04.04_uniform_distribution.py}}
\end{frame}

% Slide 12A-2: Lab Python (Bloque 2)
\begin{frame}[fragile]{Laboratorio en Python: Cuantiles y Método de Inversión}
  \scriptsize
  Verificación de cuantiles, Monte Carlo con $N=100{,}000$ y método de inversión para Exponential:
  \vspace{0.02cm}
  {\lstset{basicstyle=\fontsize{4pt}{5pt}\selectfont\ttfamily}
  \lstinputlisting[firstline=48, lastline=78]{../../code/04_variables_aleatorias_continuas/04.04_uniform_distribution.py}}
\end{frame}

% Slide 12B: Lab Python (Bloque 3)
\begin{frame}[fragile]{Laboratorio en Python: Máxima Entropía y Orden Estadísticas}
  \scriptsize
  Verificación de maximal entropía, análisis de orden estadísticas y prueba de Kolmogorov-Smirnov:
  \vspace{0.02cm}
  {\lstset{basicstyle=\fontsize{4.5pt}{5.5pt}\selectfont\ttfamily}
  \lstinputlisting[firstline=85, lastline=108]{../../code/04_variables_aleatorias_continuas/04.04_uniform_distribution.py}}
\end{frame}

% Slide 12C: Salida Terminal
\begin{frame}[fragile]{Laboratorio en Python: Salida en Terminal}
  \scriptsize
  Salida estándar generada al ejecutar el script del laboratorio en Python:
  \vspace{0.02cm}
  \begin{block}{Salida Estándar en Consola}
    \fontsize{4pt}{5pt}\selectfont
    \begin{verbatim}
=== Block 1: Uniform PDF & CDF Validation ===
Uniform(2,8) PDF integral: 1.00000000
CDF F(x) all match with scipy (diff = 0)
Quantiles q_p = a + p(b-a): all match

=== Block 2: Quantiles & Monte Carlo Generation ===
U(0,1) quantiles: q_p = p (trivial)
Monte Carlo (N=100,000): mean=0.4995 | var=0.0831
Inverse transform Exp(lambda=2): mean=0.5030 | KS=0.0046

=== Block 3: Maximum Entropy & Order Statistics ===
Entropy Uniform(0,1) = 0.0 nats (theoretical)
Entropy Beta(2,2) = -0.125 nats (< Uniform)
Order stats: E[X_(1)]=1/11, E[X_(5)]=5/11, E[X_(10)]=10/11
KS test on uniform samples: D=0.0037, p=0.1382 (fail to reject)
KS test on exponential samples: D=0.3675, p=0.0 (reject)
    \end{verbatim}
  \end{block}
\end{frame}

% Slide 13A: Ejercicio Nivel 1 (Enunciado)

% Slide 13B: Ejercicio Nivel 1 (Resolución)

% Slide 14A: Ejercicio Nivel 2 (Enunciado)

% Slide 14B: Ejercicio Nivel 2 (Resolución)

% Slide 15A: Ejercicio Nivel 3 (Enunciado)

% Slide 15B: Ejercicio Nivel 3 (Resolución)

% Slide 16A: Ejercicio Nivel 4 (Enunciado)

% Slide 16B: Ejercicio Nivel 4 (Resolución)

% Slide 17: Cierre
\begin{frame}{Cierre de Sección y Síntesis sobre la Distribución Uniforme}
  \begin{reflexion}[La Uniforme como Distribución Fundamental]
    La distribución Uniforme Continua $U(a, b)$ es la más simple de las distribuciones continuas con soporte acotado. Su rol trasciende la descripción de variables uniformes: es la base de la simulación Monte Carlo (método de inversión), la prueba no paramétrica de Kolmogorov-Smirnov, y el prior no informativo en inferencia bayesiana (por maximalidad de entropía).
  \end{reflexion}

  \vspace{0.2cm}
  \pause
  \begin{block}{Lecciones Clave de la Distribución Uniforme}
    \begin{itemize}\itemsep=0.08cm
      \item \textbf{PDF constante:} $f(x) = 1/(b-a)$ modela equiprobabilidad continua.
      \item \textbf{CDF lineal:} $F(x) = (x-a)/(b-a)$ es trivialmente invertible: $F^{-1}(p) = a + p(b-a)$.
      \item \textbf{Máxima entropía:} $U(a, b)$ es la distribución menos informativa con soporte en $[a, b]$.
    \end{itemize}
  \end{block}
\end{frame}

% Slide 18: Perspectiva Modular
\begin{frame}{Perspectiva Modular --- El Próximo Reto}
  \scriptsize
  \begin{retoclase}[Para la próxima sesión: Distribución Exponencial y Procesos sin Memoria]
    La Uniforme es la distribución continua más simple, pero la \emph{Exponencial} es quizás la más importante en aplicaciones: modela el tiempo entre eventos en un proceso de Poisson, tiene la propiedad de falta de memoria, y es la base de las distribuciones Gamma, Erlang, y Chi-cuadrada. La CDF es $F(x) = 1 - e^{-\lambda x}$, y su MLE es simplemente el inverso de la media muestral.
  \end{retoclase}

  \vspace{0.08cm}
  \pause
  \begin{alertblock}{Preparación Recomendada}
    \begin{itemize}\itemsep=0.06cm
      \item Repasa el concepto de proceso de Poisson y la conexión entre eventos Poisson y tiempos de espera exponenciales.
      \item Reflexiona sobre cómo la propiedad de falta de memoria caracteriza únicamente a la distribución exponencial.
    \end{itemize}
  \end{alertblock}
\end{frame}

\end{document}
